\documentclass[11pt]{article}
\usepackage[T1]{fontenc}
\usepackage{lmodern,amsmath,amssymb,amsthm}
\usepackage[margin=1in]{geometry}
\usepackage[colorlinks=true,linkcolor=blue,citecolor=blue,urlcolor=blue]{hyperref}
\usepackage{microtype}
\newtheorem{theorem}{Theorem}
\newtheorem{lemma}[theorem]{Lemma}
\newtheorem{corollary}[theorem]{Corollary}
\newcommand{\R}{\mathbb R}
\newcommand{\U}{\mathcal U_\rho}
\newcommand{\Y}{\mathcal Y}
\newcommand{\norm}[1]{\lVert #1\rVert}
\newcommand{\ip}[2]{\langle #1,#2\rangle}
\title{Failure of strong quadratic growth in unregularized\newline sparse control of a critical wave equation}
\author{Alec Kriebel\\\small ORCID: \href{https://orcid.org/0009-0001-9320-500X}{0009-0001-9320-500X}}
\date{October 14, 2026 \quad (version 1.0)}
\begin{document}
\maketitle
\begin{abstract}
We give an explicit counterexample to retaining ordinary critical-cone
strict positivity and strong control-$L^2$ quadratic growth after removing
quadratic control regularization from the energy-critical quintic-wave
problem of Kunisch and Meinlschmidt. On the three-dimensional unit ball,
zero is the unique global optimizer for a sufficiently small positive
pointwise control cap. Its ordinary critical cone is $\{0\}$, so the strict
condition is vacuous, but feasible controls converging in $L^\infty$ have
objective gap $o(\norm{u}_{L^2}^2)$. Nearby radial terminal targets admit
global optimizers, and every such optimizer has $L^2$ displacement divided
by target displacement tending to infinity. Local $L^1$ quadratic growth
survives. The construction realizes the established concentration and
noncoercivity mechanism in the actual nonlinear source model, using
classical wave estimates.
\end{abstract}

\section{Question, conventions and precise result}

Kunisch--Meinlschmidt's wave-control analysis \cite{KM} proves strong
control-$L^2$ quadratic growth from strict positivity on the ordinary
critical cone when the quadratic control coefficient $\beta_2$ is positive.
The Oberwolfach contribution \cite{OWR}, printed pp.~438--440, places the
$\beta_2=0$ question immediately after this implication. We answer the
question of retaining that implication unchanged. Adapted sufficient
conditions, enlarged cones and weaker growth metrics remain distinct
questions. The counterexample uses the zero critical cone explicitly.

Let $B$ be the unit ball in $\R^3$, $H=L^2(B)$ and $T=1$. All controls and
states are real. The state is the unique energy/Strichartz mild solution of
\begin{equation}\label{pde}
 y_{tt}-\Delta y+y^5=u,\qquad y|_{\partial B}=0,\qquad
 (y(0),y_t(0))=(0,0).
\end{equation}
Put $H^{-1}=H^{-1}(B)$, with the dual norm of $H_0^1(B)$, and define
\begin{align}\label{objective}
 J_{d,\nu}(u)&=\tfrac12\norm{y_u(1)-d}_H^2
       +\tfrac\nu2\norm{(y_u)_t(1)}_{H^{-1}}^2\notag\\
   &\quad+\tfrac14\norm{y_u}_{L^4(0,1;L^{12}(B))}^4
       +\norm{u}_{L^1(0,1;H)},\qquad 0\le\nu<\infty,\\
 \U&=\{u\in L^2(0,1;H):\norm{u(t)}_H\le\rho\text{ a.e.}\}.
 \notag
\end{align}
Thus $\gamma=\beta_1=1$ and $\beta_2=0$. The velocity term is the optional
term in \cite{OWR}; $\nu=0$ gives the objective in \cite{KM}. Throughout,
$\norm{u}_1$ and $\norm{u}_2$ denote time-space $L^1(H)$ and $L^2(H)$ norms;
$\norm{u}_\infty$ denotes $L^\infty(0,1;H)$.
The normalized radial eigenfunction is
\begin{equation}\label{mode}
 e(x)=\frac{\sin(\pi|x|)}{\sqrt{2\pi}\,|x|},\qquad
 e(0)=\frac\pi{\sqrt{2\pi}},\qquad
 \norm e_H=1,\quad -\Delta e=\pi^2e.
\end{equation}

\begin{theorem}\label{main}
There is $\rho_*>0$ such that, for every $0<\rho\le\rho_*$ and every fixed
finite $\nu\ge0$, the problem \eqref{objective} with $d=\pi e$ has the
following properties.
\begin{enumerate}
\item Zero is its unique global optimizer. With $A=\norm u_1$,
\begin{equation}\label{cubic}
 J_{\pi e,\nu}(u)-J_{\pi e,\nu}(0)\ge\frac{A^3}{12\rho^2}
 \quad (u\in\U).
\end{equation}
\item The ordinary critical cone at zero is $C(0)=\{0\}$. The source's
substitute second-variation strict condition on $C(0)\setminus\{0\}$ holds
vacuously.
\item The feasible controls $u_h=h e\,\mathbf1_{|t-1/2|<h}$, for
$0<h<\min(\rho,1/2)$, satisfy
\begin{equation}\label{pulseclaim}
 \norm{u_h}_\infty=h,\quad \norm{u_h}_2^2=2h^3,\quad
 J_{\pi e,\nu}(u_h)-J_{\pi e,\nu}(0)
  =\left(\frac{\pi^2}{3}+\frac2{\pi^2}\right)h^4+O(h^6).
\end{equation}
Consequently no positive control-$L^2$ quadratic-growth constant exists
in either an $L^2$ or an $L^\infty$ neighborhood of zero.
\item For $d_\varepsilon=(\pi+\varepsilon)e$, global optimizers exist and
every choice of global optimizers satisfies
\begin{equation}\label{instability}
 \frac{\norm{u_\varepsilon}_2}{\varepsilon}\longrightarrow\infty
 \qquad(\varepsilon\downarrow0).
\end{equation}
\end{enumerate}
The cap $\rho_*$ and the constants in \eqref{cubic} can be chosen
independently of $\nu$. Small neighborhoods in the radial perturbation
assertion may depend on the fixed $\nu$.
\end{theorem}

The distinction between exact-cone positivity and coercivity against
normalized weakly null sequences is established in second-order control
theory \cite{CT12,CHW,CW,BW}. The Dunn example, as reproduced in
\cite[Example~4.16]{CT15}, already illustrates exact-cone insufficiency.
Group-sparse threshold calculus is classical \cite{HSW}. Our contribution
is the explicit realization with a genuine nonlinear-wave optimizer and
the quantitative assertions above. An extensive bounded primary-source
audit located no earlier covering realization; it does not certify
worldwide firstness. All analytic estimates used below are standard
inputs from \cite{KM}.

\section{Uniform small-input estimates and global optimality}

Use the compatible state space
\[
 \Y=\{y\in C([0,1];H_0^1(B)):y_t\in C([0,1];H),\
                    y\in L^4(0,1;L^{12}(B))\},
\]
with the sum norm and the distributional derivative compatibility.
Let $L u=z_u$ denote the zero-data linear-wave solution. The bounded-domain
energy/Strichartz estimate in \cite[Lemma~2.9]{KM} gives
$\norm{Lu}_{\Y}\le K\norm u_1$. Sobolev embedding and interpolation give
\begin{equation}\label{interp}
 \norm y_{L^5L^{10}}^5\le
 C\norm y_{L^4L^{12}}^4\norm y_{L^\infty H_0^1}.
\end{equation}

\begin{lemma}\label{small}
There are $a_0,E,R>0$ such that, for every $u$ with $A=\norm u_1\le a_0$,
\begin{equation}\label{remainder}
 \norm{y_u}_{\Y}\le E A,\qquad
 \norm{y_u-z_u}_{\Y}\le R A^5.
\end{equation}
The constants are independent of the terminal target and $\nu$.
\end{lemma}
\begin{proof}
On the closed compatible zero-data subspace of $\Y$, consider
$y\mapsto L(u-y^5)$. By \eqref{interp}, the quintic has norm at most
$M\norm y_\Y^5$ in $L^1(H)$ and difference norm at most
$M(\norm y_\Y+\norm w_\Y)^4\norm{y-w}_\Y$.
For $A\le a_0$ sufficiently small this map preserves the radius-$2KA$
ball and is a contraction there. The resulting mild solution agrees with
the unique energy/Strichartz state in \cite{KM}. Its norm is at most
$2KA$, and $y_u-z_u=-L(y_u^5)$ proves the second bound. Enlarge $E,R$
as necessary to bound terminal position in $H$ and terminal velocity in
$H^{-1}$. The compatible residual $y+L(y^5)-Lu$ has derivative the identity
in $y$ at zero. Its quintic part is a bounded polynomial by
\eqref{interp}, so the local implicit-function theorem also gives a
smooth state map with $S'(0)=L$ and $S''(0)=0$.
This direct argument does not require an inverse onto arbitrary
continuous position/velocity pairs.
\end{proof}

For $u\in\U$, write $f(t)=\norm{u(t)}_H$, $A=\int f$ and $r=\norm u_2$.
Then
\begin{equation}\label{cap}
 0\le A\le\rho,\qquad A\le r,\qquad r^2\le\rho A.
\end{equation}
Projection of the linear state onto \eqref{mode} gives
\begin{equation}\label{endpoint}
 \ip e{z_u(1)}=\frac1\pi\int_0^1\sin(\pi t)\ip e{u(t)}\,dt.
\end{equation}
This is only a projection identity; the nonlinear state need not stay in
the span of $e$. Expanding the terminal square, discarding nonnegative
terms, and using \eqref{remainder} yield
\begin{equation}\label{weighted}
 \Delta(u):=J_{\pi e,\nu}(u)-J_{\pi e,\nu}(0)
 \ge\int_0^1(1-\sin\pi t)f(t)\,dt-C_0 A^5,
\end{equation}
where $C_0$ is fixed independently of $\nu$.

Since $1-\sin(\pi t)\ge2(t-1/2)^2$, a same-mass comparison gives
\begin{equation}\label{moment}
 \int_0^1(t-1/2)^2 f(t)\,dt\ge\frac{A^3}{12\rho^2}.
\end{equation}
Indeed, put $b=A/(2\rho)\le1/2$ and
$f_*=\rho\mathbf1_{|t-1/2|<b}$. The product
$((t-1/2)^2-b^2)(f-f_*)$ is nonnegative on both the centered interval
and its complement. Its integral proves \eqref{moment}, because
$\int(f-f_*)=0$. The cases $A=0$ and $A=\rho$ are included.
Choose $\rho_*\le a_0$ with $C_0\rho_*^4\le1/12$. Equations
\eqref{weighted}--\eqref{moment} then give
\[
 \Delta(u)\ge\frac{A^3}{6\rho^2}-C_0A^5
              \ge\frac{A^3}{12\rho^2}.
\]
This proves global uniqueness in Theorem~\ref{main}. No numerical value
for an unknown PDE constant is asserted.

\section{Ordinary cone and loss of strong growth}

In the source convention, the ordinary critical cone consists of tangent
directions $v$ with zero first directional derivative of the whole cost.
At zero, bounded truncations of any $v\in L^1(H)$ or $L^2(H)$ are radially
feasible after sufficiently small scaling. They converge in the respective
norm, so the closed tangent cone is the whole space. Pointwise constraint
slack alone would not imply norm interiority in these two spaces.
The adjoint and sparse subgradient are
\[
 p(t)=-\sin(\pi t)e,\qquad \lambda(t)=\sin(\pi t)e,
 \qquad\norm{\lambda(t)}_H\le1,
\]
with zero constraint multiplier. Direct differentiation gives
\begin{equation}\label{first}
 DJ_{\pi e,\nu}(0;v)
  =\int_0^1\big(\norm{v(t)}_H-\sin(\pi t)\ip e{v(t)}\big)\,dt
  \ge\int_0^1(1-\sin\pi t)\norm{v(t)}_H\,dt.
\end{equation}
The right-hand side is positive for every nonzero Lebesgue equivalence
class: its weight vanishes only at $t=1/2$. Hence $C(0)=\{0\}$.
The source explicitly sets the substitute sparse $j''(0;v^2)$ to zero;
this is not a classical second derivative of a norm. Its strict condition
on nonzero critical directions is satisfied vacuously.

For the pulse in Theorem~\ref{main}, direct linear-wave integration gives
\begin{equation}\label{pulse}
 A_h=2h^2,\quad r_h^2=2h^3,\quad
 z_{u_h}(1)=\frac{2h\sin(\pi h)}{\pi^2}e,\quad
 (z_{u_h})_t(1)=0.
\end{equation}
Its linear tracking-plus-sparse gap is
\[
 2h^2-\frac{2h\sin(\pi h)}\pi
       +\frac{2h^2\sin^2(\pi h)}{\pi^4}
 =\left(\frac{\pi^2}{3}+\frac2{\pi^2}\right)h^4+O(h^6).
\]
Lemma~\ref{small} bounds the nonlinear terminal error by $O(h^{10})$,
the quartic state cost by $O(h^8)$, and the velocity cost by
$O(\nu h^{20})$. This proves \eqref{pulseclaim}. Dividing by $2h^3$
proves the claimed failure of growth, even with $\norm{u_h}_\infty\to0$.

For comparison with the established NDC framework \cite{CW,BW}, set
$v_h=u_h/r_h$. These unit-norm $L^2(H)$ directions are weakly null:
their pairing with an $L^2$ test is bounded by the test's norm on the
shrinking interval. With $F$ the smooth costs and
$G=\norm{\cdot}_1+I_{\U}$, the normalized second-order expression is
\begin{align*}
 \frac{G(u_h)-G(0)+F'(0)u_h}{r_h^2}
   +\tfrac12 F''(0)[v_h,v_h]
 &=\frac1h-\frac{\sin(\pi h)}{\pi h^2}
                      +\frac{\sin^2(\pi h)}{\pi^4h}\\
 &=\left(\frac{\pi^2}{6}+\frac1{\pi^2}\right)h+O(h^3)\to0.
\end{align*}
Thus exact-cone positivity misses this familiar weak-null obstruction.
The direct pulse argument already handles both source localizations.

\section{Existence and terminal-target perturbations}

We record existence for the optional velocity term as well. The set $\U$
is closed, convex and bounded in $L^2(H)$, hence weakly compact. Along a
minimizing sequence, Lemma~\ref{small} bounds the states uniformly in
energy and $L^4L^{12}$. Compactness of $H_0^1\hookrightarrow H$ and the
bound on $y_t$ imply strong compactness of positions in $C([0,1];H)$.
Moreover $y_{tt}=\Delta y-y^5+u$ is bounded in $L^\infty(H^{-1})$:
$y\in L^\infty L^6$, so $y^5\in L^\infty L^{6/5}\hookrightarrow
L^\infty H^{-1}$. Compactness of $H\hookrightarrow H^{-1}$ then gives
strong compactness of velocities in $C([0,1];H^{-1})$.

A subsequence of states converges almost everywhere. The bounded
$L^{6/5}$ quintics converge weakly to the quintic of the limit; the limit
also retains the weak energy and $L^4L^{12}$ bounds. Consequently its
quintic lies in $L^1(H)$ by \eqref{interp}. Passing to the distributional
equation and initial traces identifies a zero-data weak solution.
Subtract the linear Duhamel solution forced by $u-y^5$. The remaining
homogeneous Dirichlet wave has zero initial data; projection onto a
complete Dirichlet eigenbasis proves it is zero. Thus the limit is the
actual mild state, by uniqueness. Endpoint costs are continuous under the
two strong convergences, and the sparse and quartic norms are weakly lower
semicontinuous. The direct method therefore gives global optimizers for
every $d\in H$ and every finite $\nu\ge0$ in this capped small-input setting.

Fix $a=\rho/2$ and use the feasible trial
\[
 \widehat u_\varepsilon=a e\,
 \mathbf1_{|t-1/2|<\pi\varepsilon/(2a)}.
\]
For sufficiently small $\varepsilon$, its mass is $\pi\varepsilon$,
its linear endpoint is $\varepsilon e+O(\varepsilon^3)$ and its linear
terminal velocity is zero. Expansion as above gives, for fixed $\nu$,
\begin{equation}\label{negative}
 J_{d_\varepsilon,\nu}(\widehat u_\varepsilon)
       -J_{d_\varepsilon,\nu}(0)
    =-\tfrac12\varepsilon^2+O(\varepsilon^3)
    \le-\tfrac14\varepsilon^2.
\end{equation}
For all feasible $u$, discarding nonnegative costs gives instead
\begin{equation}\label{pertlower}
 J_{d_\varepsilon,\nu}(u)-J_{d_\varepsilon,\nu}(0)
 \ge-\norm{g_\varepsilon}_{L^2(0,1)}r-C_1 A^5,
 \quad
 g_\varepsilon(t)=\big[(1+\varepsilon/\pi)\sin(\pi t)-1\big]_+.
\end{equation}
For small $\varepsilon$,
\begin{equation}\label{dualbound}
 g_\varepsilon(t)\le[\varepsilon/\pi-2(t-1/2)^2]_+,
 \qquad
 \norm{g_\varepsilon}_{L^2}
       \le 2^{1/4}\pi^{-5/4}\varepsilon^{5/4}.
\end{equation}
If an optimizer had $r\le K\varepsilon$, then $A\le r$ and
\eqref{pertlower} would be $-o(\varepsilon^2)$, contradicting
\eqref{negative}. For every fixed $K$ this is impossible for sufficiently
small $\varepsilon$, proving \eqref{instability} for every optimizer.

\section{A surviving weaker-norm estimate}

\begin{corollary}\label{weaker}
For the same example there is $\delta>0$ with
\begin{equation}\label{L1growth}
 \Delta(u)\ge\frac{\norm u_1^2}{4\pi^2}
 \quad (u\in\U,\ \norm u_1\le\delta).
\end{equation}
For any target $\pi e+w$, with $\varepsilon=\norm w_H$ sufficiently small,
every global optimizer satisfies
\begin{equation}\label{upperrates}
 \norm u_1\le4\pi^2E\varepsilon,\qquad
 \norm u_2\le2\pi\sqrt{\rho E}\sqrt\varepsilon.
\end{equation}
The constants and this upper-bound neighborhood can be chosen uniformly
for all finite $\nu\ge0$. On $w=\varepsilon e$, for each fixed $\nu$,
\begin{equation}\label{radialrates}
 \norm{u_\varepsilon}_1=\Theta(\varepsilon),\qquad
 c\varepsilon^{3/4}\le\norm{u_\varepsilon}_2\le C\sqrt\varepsilon
\end{equation}
for every global optimizer and sufficiently small $\varepsilon$.
The lower-bound neighborhood in \eqref{radialrates} may depend on $\nu$.
\end{corollary}
\begin{proof}
Set $x=\ip e{y_u(1)}\le b=A/\pi+RA^5$. The full tracking square gives
$\Delta(u)\ge q_A(x)=A-\pi x+x^2/2$. Choose $\delta$ so that $b\le\pi/2$
and $\pi R\delta^3\le1/(4\pi^2)$. Since $q_A$ decreases for $x<\pi$,
\[
 \Delta(u)\ge q_A(b)
 =\frac{A^2}{2\pi^2}-\pi RA^5+\frac R\pi A^6+\frac{R^2}2A^{10}
 \ge\frac{A^2}{4\pi^2}.
\]
For a perturbed global optimizer, comparison with zero gives
$\Delta(u)\le\ip w{y_u(1)}\le E\varepsilon A$. First \eqref{cubic}
implies $A\le\sqrt{12\rho^2E\varepsilon}$, placing it in the $\delta$
neighborhood. Then \eqref{L1growth} yields the first upper bound, and
$r^2\le\rho A$ gives the second. Every step discards the nonnegative
velocity cost, so these upper conclusions are uniform in $\nu$.

On the radial ray, the upper $A=O(\varepsilon)$ makes the last term in
\eqref{pertlower} $O(\varepsilon^5)$. Combining \eqref{negative} with
\eqref{dualbound} gives $r\ge c\varepsilon^{3/4}$. Also
$\norm{g_\varepsilon}_\infty\le\varepsilon/\pi$, so the same comparison
gives $A\ge c'\varepsilon$. This proves \eqref{radialrates}, with the
smallness threshold inherited from the fixed-$\nu$ trial comparison.
\end{proof}

These are estimates relative to the fixed zero optimizer. They assert no
pairwise Lipschitz selection, perturbed optimizer uniqueness,
$L^\infty$ convergence, or sharpness of the exponent bracket.
Weaker-norm and strengthened-cone theories are developed in
\cite{WW,CHW,CW,BW}; the elementary estimate here uses the terminal
tracking square and the cap. For example, taking target zero in the same
model gives $J_{0,\nu}(u)\ge\norm u_1\ge\norm u_2^2/\rho$.
Thus the theorem is a specific obstruction to the unchanged implication,
not universal failure of all unregularized quadratic growth.

\paragraph{Verification and review status.}
The accompanying package contains this source, exact algebra/control
checkers, reference receipts and independently authored proof/priority
records. Finite checks do not numerically certify a nonlinear PDE solution,
an unknown PDE constant, or historical novelty. AI tools were used
extensively for construction, derivation, independent adversarial audits,
computation and writing. This is an unrefereed research note; no external
human peer review is claimed. The author takes responsibility for the
stated mathematics. The package identifies the bounded priority audit and
its access limitations.

\begin{thebibliography}{10}
\bibitem{KM} K. Kunisch and H. Meinlschmidt,
Optimal control of an energy-critical semilinear wave equation in 3D with
spatially integrated control constraints, \emph{J. Math. Pures Appl.}
138 (2020), 46--87. \href{https://doi.org/10.1016/j.matpur.2020.03.006}{doi:10.1016/j.matpur.2020.03.006}.
Proof version inspected: \href{https://arxiv.org/pdf/1907.02744v1}{arXiv:1907.02744v1} (2019).
\bibitem{OWR} H. Meinlschmidt, joint work with K. Kunisch,
Optimal Control of a Semilinear Critical Wave Equation, in
\emph{Challenges in Optimization with Complex PDE-Systems}, Oberwolfach
Report 9/2021, 438--440.
\href{https://doi.org/10.4171/OWR/2021/9}{doi:10.4171/OWR/2021/9}.
\bibitem{HSW} R. Herzog, G. Stadler and G. Wachsmuth,
Directional Sparsity in Optimal Control of Partial Differential Equations,
\emph{SIAM J. Control Optim.} 50 (2012), 943--963.
\href{https://doi.org/10.1137/100815037}{doi:10.1137/100815037}.
\bibitem{CT12} E. Casas and F. Tr\"oltzsch,
Second Order Analysis for Optimal Control Problems: Improving Results
Expected From Abstract Theory, \emph{SIAM J. Optim.} 22 (2012), 261--279.
\href{https://doi.org/10.1137/110840406}{doi:10.1137/110840406}.
\bibitem{CT15} E. Casas and F. Tr\"oltzsch,
Second Order Optimality Conditions and Their Role in PDE Control,
\emph{Jahresber. Dtsch. Math.-Ver.} 117 (2015), 3--44.
\href{https://doi.org/10.1365/s13291-014-0109-3}{doi:10.1365/s13291-014-0109-3}.
\bibitem{CHW} E. Casas, R. Herzog and G. Wachsmuth,
Analysis of Spatio-Temporally Sparse Optimal Control Problems of Semilinear
Parabolic Equations, \emph{ESAIM COCV} 23 (2017), 263--295.
\href{https://doi.org/10.1051/cocv/2015048}{doi:10.1051/cocv/2015048}.
\bibitem{CW} C. Christof and G. Wachsmuth,
No-Gap Second-Order Conditions via a Directional Curvature Functional,
\emph{SIAM J. Optim.} 28 (2018), 2097--2130.
\href{https://doi.org/10.1137/17M1140418}{doi:10.1137/17M1140418}.
\bibitem{BW} N. Borchard and G. Wachsmuth,
Second-order conditions for spatio-temporally sparse optimal control via
second subderivatives, \emph{J. Nonsmooth Anal. Optim.} 5 (2024), 12604.
\href{https://doi.org/10.46298/jnsao-2024-12604}{doi:10.46298/jnsao-2024-12604}.
Original preprint: arXiv:2311.14538v1 (2023).
\bibitem{WW} D. Wachsmuth and G. Wachsmuth,
Second-order conditions for non-uniformly convex integrands: quadratic
growth in $L^1$, \emph{J. Nonsmooth Anal. Optim.} 3 (2022), 8733.
\href{https://doi.org/10.46298/jnsao-2022-8733}{doi:10.46298/jnsao-2022-8733}.
\end{thebibliography}
\end{document}
